\chapter{公式清单}\label{appendix:公式清单}

    \section{pH 与溶液酸碱性的计算}\label{app:公式-pH}

        \begin{itemize}
            \item \(\cemh{pH} = -\lg{}\cemh{[H+]}\)，\(\cemh{pOH} = -\lg{}\cemh{[OH-]}\)，适用范围 \numrange{0}{14}。
            \item 水的离子积：\(K_\text{W} = \cemh{[H+]} \cdot \cemh{[OH-]}\)（常温 \(\num{1E-14}\)），\(\cemh{[H+]} = \dfrac{K_\text{W}}{\cemh{[OH-]}}\)。
            \item 强酸、强碱每稀释 10 倍，\cemh{[H+]}、\cemh{[OH-]} 变化 10 倍，pH 改变 1；\textcolor{red}{酸不会稀释成碱、碱不会稀释成酸，pH 最终逼近 7}。
            \item 酸碱混合先判断过量：酸过量算 \cemh{[H+]}（再求 pH），碱过量算 \cemh{[OH-]}（再求 pH）。
            \item 若反应不涉及 \cemh{H+}、\cemh{[OH-]}，相当于稀释。
            \item 非常温下先由中性溶液的 pH 算出该温度下 \(K_\text{W}\)，其余计算同前。
        \end{itemize}

    \section{弱电解质的电离平衡}\label{app:公式-电离平衡}

        \subsubsection{\texorpdfstring{电离度 \(\alpha\)}{电离度}}

            \[
                \alpha = \frac{\text{已电离的分子数}}{\text{分子总数}} \times \SI{100}{\percent} = \frac{\cemh{[H+]}}{c} \times \SI{100}{\percent}
            \]

        \subsubsection{\texorpdfstring{电离常数 \(K_\text{a}\)}{电离常数}}

            对于 \cemh{HA <=> H+ + A-}：\[
                K_\text{a} = \frac{\cemh{[H+]} \cdot \cemh{[A-]}}{\cemh{[HA]}}
            \]

        \subsubsection{\texorpdfstring{近似计算（\(\frac{c}{K_\text{a}} \ge 400\)）}{近似计算}}

            \[
                \cemh{[H+]} = \sqrt{cK_\text{a}} \qquad \alpha = \sqrt{\frac{K_\text{a}}{c}}
            \]

            \textcolor{red}{弱电解质浓度越大，电离度越小（越稀越电离）。}

    \section{盐类水解}\label{app:公式-水解}

        \subsubsection{\texorpdfstring{水解度 \(h\)}{水解度}}

            \[
                h = \frac{\text{已水解的离子数}}{\text{离子总数}} \times \SI{100}{\percent} = \frac{\cemh{[OH-]}}{c} \times \SI{100}{\percent}
            \]

        \subsubsection{\texorpdfstring{水解常数 \(K_\text{h}\)}{水解常数}}

            对于 \cemh{A- + H2O <=> HA + OH-}：\[
                K_\text{h} = \frac{\cemh{[HA]} \cdot \cemh{[OH-]}}{\cemh{[A-]}} = \frac{K_\text{W}}{K_\text{a}}
            \]

            近似计算：\[
                \cemh{[OH-]} = \sqrt{cK_\text{h}} \qquad h = \sqrt{\frac{K_\text{h}}{c}}
            \]

            \textcolor{red}{盐浓度越大，水解度越小（越稀越水解）。}

            \textcolor{blue}{多元弱酸根分步水解：\cemh{CO3^{2-}} 的 \(K_\text{h1} = \dfrac{K_\text{W}}{K_\text{a2}}\)、\(K_\text{h2} = \dfrac{K_\text{W}}{K_\text{a1}}\)，且第一步远大于第二步，故碱性正盐 \(>\) 酸式盐。}

    \section{沉淀溶解平衡}\label{app:公式-溶度积}

        \subsubsection{\texorpdfstring{溶度积 \(K_\text{sp}\)}{溶度积}}

            对于 \cemh{A_mB_n(s) <=> mA^{n+} + nB^{m-}}：\[
                K_\text{sp} = \cemh{[A^{n+}]}^m \cdot \cemh{[B^{m-}]}^n
            \]

        \subsubsection{溶度积规则}

            \begin{itemize}
                \item \(Q < K_\text{sp}\)：无沉淀析出。
                \item \(Q = K_\text{sp}\)：沉淀溶解平衡状态。
                \item \(Q > K_\text{sp}\)：有沉淀析出，直至达到平衡。
            \end{itemize}

        \subsubsection{\texorpdfstring{溶解度 \(s\) 与 \(K_\text{sp}\)}{溶解度与溶度积}}

            \begin{datas}
                \item \cemh{AB} 型（如 \cemh{AgCl}、\cemh{BaSO4}）：\(s = \sqrt{K_\text{sp}}\)。
                \item \cemh{A2B} 或 \cemh{AB2} 型（如 \cemh{Ag2S}）：\(s = \sqrt[3]{\dfrac{K_\text{sp}}{4}}\)。
            \end{datas}

            \textcolor{blue}{同类型难溶电解质，\(K_\text{sp}\) 越小溶解度越小。}

        \subsubsection{沉淀转化与沉淀溶解的平衡常数}

            沉淀转化（如 \cemh{AgCl + I- \xleq{} AgI + Cl-}）：\[
                K = \frac{K_\text{sp}(\text{前沉淀})^\text{幂次方}}{K_\text{sp}(\text{后沉淀})^\text{幂次方}}
            \]

            沉淀溶于酸生成水（如 \cemh{Mg(OH)2 + 2H+}）：\[
                K = \frac{K_\text{sp}}{K_\text{W}^2}
            \]

            沉淀溶于酸生成气体（如 \cemh{CaCO3 + 2H+}）：\[
                K = \frac{K_\text{sp}}{K_\text{a1} \cdot K_\text{a2}}
            \]
            其中 \(K_\text{a1}\)、\(K_\text{a2}\) 为所生成的弱酸（如 \cemh{H2CO3}）的各级电离常数。

            \textcolor{blue}{方法二：转化平衡常数可拆为两步骤的 \(K_1 \times K_2\) 求得。}

    \section{溶液中的三大守恒}\label{app:公式-守恒}

        \begin{itemize}
            \item \textcolor{red}{电荷守恒}：溶液中正电荷总数 \(=\) 负电荷总数（记得乘电荷数）。如 \(\cemh{Na2CO3}\) 溶液中：
            \[
                \cemh{[Na+]} + \cemh{[H+]} = \cemh{[OH-]} + \cemh{[HCO3-]} + 2\cemh{[CO3^{2-}]}
            \]
            \item \textcolor{red}{物料守恒}：元素总量守恒（微粒要找全）。如 \(\cemh{Na2CO3}\) 溶液中：
            \[
                \cemh{[Na+]} = 2\left(\cemh{[CO3^{2-}]} + \cemh{[HCO3-]} + \cemh{[H2CO3]}\right)
            \]
            \item \textcolor{red}{质子守恒}：由电荷守恒与物料守恒叠加得到。如 \(\cemh{Na2CO3}\) 溶液中：
            \[
                \cemh{[OH-]} = \cemh{[H+]} + \cemh{[HCO3-]} + 2\cemh{[H2CO3]}
            \]
        \end{itemize}

    \section{中和滴定}\label{app:公式-滴定}

        \[
            c_\text{酸} V_\text{酸} N_\text{酸} = c_\text{碱} V_\text{碱} N_\text{碱}
        \]

        \(N\) 为酸（碱）的元数（如 \cemh{H2SO4} 的 \(N = 2\)）；\(c\) 为浓度，\(V\) 为体积。

    \section{电化学}\label{app:公式-电化学}

        \begin{itemize}
            \item \textcolor{red}{电子守恒}：电极上得电子总数 \(=\) 失电子总数（配平及计算依据）。
            \item 配平依据：化合价升降总数相等；离子方程式顺序——电子守恒 \(\to\) 电荷守恒 \(\to\) 原子守恒。
            \item \textcolor{red}{极区 pH 变化看电极反应，不看总反应}（如电解 \cemh{NaCl} 阴极区生成 \cemh{OH-}，pH 增大）。
        \end{itemize}

    \section{晶体结构}\label{app:公式-晶体}

        \subsubsection{晶体密度计算}

            \[
                \rho = \frac{\text{晶胞质量}}{\text{晶胞体积}}
                = \frac{\dfrac{M}{N_\text{A}} \times \text{晶胞中粒子数}}{a^3}
            \]

            \textcolor{blue}{思路：根据晶胞结构确定各种微粒数目 \(\to\) 晶胞质量 \(\to\) 密度 \(=\) 晶胞质量 \(\div\) 晶胞体积。已知微粒半径时，由半径先求晶胞边长。}

        \subsubsection{晶胞中粒子数的计算（均摊法）}

            \[
                \text{粒子数} = \text{顶角} \times \frac{1}{8} + \text{棱上} \times \frac{1}{4} + \text{面心} \times \frac{1}{2} + \text{体心} \times 1
            \]

        \subsubsection{空间利用率}

            \begin{datas}
                \item 简单立方堆积：\SI{52.36}{\percent}
                \item 体心立方堆积：\SI{68.02}{\percent}
                \item 面心立方、六方最密堆积：\SI{74.05}{\percent}
            \end{datas}
